\scanpage{21}
\margin{en admettant la résolution des singularités ... et les conjectures de Weil\par\rule{.32em}{.3pt}\ \rule{.65em}{.3pt}\ \rule{.43em}{.3pt}}
\[
\boxed{L_E\left(\frac{1}{p^{n+\rho}t}\right)=A(t)\,\delta(E)\,(-t)^{\chi(E)}\,L_E(t)}
\]
\[
\left\{\begin{aligned}
A(t)&=L_E^\infty\left(\deleted{p}\frac{1}{p^{n+\rho}t}\right)^{-1}\\
&\deleted{\unsure{=|\simeq}}\\
\chi(E)&=\sum(-1)^i\operatorname{rang}H^i_!(\bar X,\bar F)\\
\delta(E)&=\prod_i\deleted{\unsure{(-1)^i|\varepsilon_i}}\bigl(\operatorname{dét}\operatorname{frob}_{H^i_!(\bar X,\bar F)}\bigr)^{(-1)^i}
\end{aligned}\right.
\]

\uline{NB}\quad $L_E^\infty(t)=\inserted{=L^\infty_{\mathbb R_\infty f_*(E)}(t)=}L_{\mathbb R i_*(E)|Y}(t)$ est de \deleted{\unsure{niveau|niveau $\rho$}} \inserted{type} $\leq\tau+(n-1)$ ($\tau$ étant le type de $E$) et de poids $\leq\rho+2n$

\uline{Si $X$ complet}
\[
A(t)=1
\]

On peut aussi écrire
\[
\chi(E)=\varepsilon(\mathbb R f_!(E)),\quad \deleted{\unsure{\varepsilon|=}}\delta(E)=\deleted{\unsure{F|f}\,\operatorname{dét}\unsure{\operatorname{frob}|f}}\delta(\mathbb R f_!(E))
\]
\[
\varepsilon(\mathbb R f_!(E))=\varepsilon(D\mathbb R f_!(E))\deleted{\unsure{=|\varepsilon}},\quad D\mathbb R f_!(E)=\mathbb R\unsure{f|f_*}(DE)=(\mathbb R f_*(E))(\rho+\deleted{\unsure{2|2n}}n)\inserted{[2n]}
\]
\[
\deleted{\operatorname{det}}\delta(\mathbb R f_!(E))=\delta(D(\mathbb R f_!(E))^{-1}
\]
d'où \deleted{d'où}
\[
\deleted{\unsure{\lbrace|f}}\left\{\begin{aligned}
\chi(E)&=\sum(-1)^i\operatorname{rang}\mathbb R^i f_*(E)\\
\delta(E)&=p^{+(\rho+n)\chi(E)}\delta(\mathbb R f_*(E))^{-1}=p^{(n+\rho)\chi(E)}\prod_i\deleted{\unsure{(-1)^i|\delta_i}}\bigl(\operatorname{det} f_{H^i(\bar X,\bar F)}\bigr)^{(-1)^{i+1}}
\end{aligned}\right.
\]

[Comparant, lorsque $X$ est complet, les deux valeurs obtenues de $\delta(E)$, on trouve dans ce cas
\[
\delta(E)^2=p^{(n+\rho)\chi(E)}\qquad (X\text{ complet})
\]
ce qui est bien compatible avec les conjectures de Weil sur les val. absolues des valeurs propres.... ]
% archive sheet number: 20
